% Seminar 2 - Pagina 8. Fragment pentru % Seminar 2 - Pagina 8. Fragment pentru % Seminar 2 - Pagina 8. Fragment pentru % Seminar 2 - Pagina 8. Fragment pentru \input{seminars/pages/seminar2-pagina8-ro}.
% Fara preambul si fara intrerupere fortata de pagina.
% NOTE EDITORIALE:
% - Primele formule continua inductia pentru principiul includerii si
%   excluderii din pagina 7. Semnele + din expresia intermediara cu
%   (A_i \cap A_{n+1}) sunt redate ca in manuscris, desi semnul necesar
%   pentru formula finala ar fi minus; nu s-a facut corectarea tacita.
% - In problema 11, manuscrisul foloseste A_i' la definitia multimii
%   de functii care omit b_i, dar omite uneori prima in formulele
%   pentru reuniune si intersectie. In LaTeX se pastreaza A_i' in
%   formulele de numarare pentru a indica aceeasi familie de multimi.
% - La sfarsitul ultimei sume apare un termen taiat / rescris; se reda
%   ultima forma lizibila +(-1)^{n-1} 0, fara a reconstitui termenul taiat.
% - Schemele marginale reprezinta doar exemple de corespondente intre
%   multimi finite; afirmatia explicita s(2,1)=1 este pastrata.

\subsection*{Seminar 2 - Pagina 8}

\noindent\textit{(Continuarea demonstraţiei principiului includerii şi excluderii.)}
\[
\begin{aligned}
&\cdots+(-1)^{n-1}\bigl|A_1\cap\cdots\cap A_n\bigr|\\
&\quad+\bigl|(A_1\cap A_{n+1})\cup(A_2\cap A_{n+1})
    \cup\cdots\cup(A_n\cap A_{n+1})\bigr|\\
&=\sum_{i=1}^{n+1}|A_i|
  -\sum_{1\leq i<j\leq n}|A_i\cap A_j|
  +\cdots+(-1)^{n-1}|A_1\cap\cdots\cap A_n|\\
&\quad+\sum_{i=1}^{n}|A_i\cap A_{n+1}|
  -\sum_{1\leq i<j\leq n}|A_i\cap A_j\cap A_{n+1}|
  +\cdots\\
&\quad+(-1)^{n-1}|A_1\cap A_2\cap\cdots\cap A_n\cap A_{n+1}|.
\end{aligned}
\]
\[
\begin{aligned}
&=\sum_{i=1}^{n+1}|A_i|
  -\sum_{1\leq i<j\leq n+1}|A_i\cap A_j|\\
&\quad+\sum_{1\leq i<j<k\leq n+1}|A_i\cap A_j\cap A_k|
  -\cdots\\
&\quad+(-1)^n|A_1\cap A_2\cap\cdots\cap A_n\cap A_{n+1}|.
\end{aligned}
\]

\bigskip
\noindent\fbox{11}\quad $|M|=m$, $|N|=n$.

\noindent Definim
\[
 s(m,n)=\bigl|\{f:M\longrightarrow N\mid f\text{ surjectivă}\}\bigr|.
\]
\noindent Calculaţi $s(m,n)$. \hfill $s(2,1)=1$.

\noindent\underline{Sol.}
\[
 s(m,n)=0\qquad\text{când }n>m.
\]
\noindent Dacă $n=m$, atunci
\[
 s(m,n)=n!.
\]
\noindent Dacă $n<m$, atunci
\[
 |\operatorname{Hom}(M,N)|=n^m.
\]
\noindent Fie $P=\{f:M\longrightarrow N\mid f\text{ nu este surjectivă}\}$.
Atunci
\[
 s(m,n)=n^m-|P|.
\]
\noindent Scriem
\[
 M=\{a_1,a_2,\ldots,a_m\},\qquad
 N=\{b_1,b_2,\ldots,b_n\}.
\]
\noindent Pentru $i=1,\ldots,n$ notăm
\[
 A_i'=\{f:M\longrightarrow N\mid b_i\notin\operatorname{Im}f\}.
\]
\noindent Rezultă
\[
 |P|=\bigl|A_1'\cup A_2'\cup\cdots\cup A_n'\bigr|,
\]
\[
 |A_i'|=(n-1)^m\quad(\forall i),\qquad
 |A_i'\cap A_j'|=(n-2)^m.
\]
\noindent Aplicând principiul includerii şi excluderii,
\[
\begin{aligned}
 |P|={}&\binom{n}{1}(n-1)^m
       -\binom{n}{2}(n-2)^m\\
      &+\binom{n}{3}(n-3)^m-\cdots
       +(-1)^{n-1}0.
\end{aligned}
\]
.
% Fara preambul si fara intrerupere fortata de pagina.
% NOTE EDITORIALE:
% - Primele formule continua inductia pentru principiul includerii si
%   excluderii din pagina 7. Semnele + din expresia intermediara cu
%   (A_i \cap A_{n+1}) sunt redate ca in manuscris, desi semnul necesar
%   pentru formula finala ar fi minus; nu s-a facut corectarea tacita.
% - In problema 11, manuscrisul foloseste A_i' la definitia multimii
%   de functii care omit b_i, dar omite uneori prima in formulele
%   pentru reuniune si intersectie. In LaTeX se pastreaza A_i' in
%   formulele de numarare pentru a indica aceeasi familie de multimi.
% - La sfarsitul ultimei sume apare un termen taiat / rescris; se reda
%   ultima forma lizibila +(-1)^{n-1} 0, fara a reconstitui termenul taiat.
% - Schemele marginale reprezinta doar exemple de corespondente intre
%   multimi finite; afirmatia explicita s(2,1)=1 este pastrata.

\subsection*{Seminar 2 - Pagina 8}

\noindent\textit{(Continuarea demonstraţiei principiului includerii şi excluderii.)}
\[
\begin{aligned}
&\cdots+(-1)^{n-1}\bigl|A_1\cap\cdots\cap A_n\bigr|\\
&\quad+\bigl|(A_1\cap A_{n+1})\cup(A_2\cap A_{n+1})
    \cup\cdots\cup(A_n\cap A_{n+1})\bigr|\\
&=\sum_{i=1}^{n+1}|A_i|
  -\sum_{1\leq i<j\leq n}|A_i\cap A_j|
  +\cdots+(-1)^{n-1}|A_1\cap\cdots\cap A_n|\\
&\quad+\sum_{i=1}^{n}|A_i\cap A_{n+1}|
  -\sum_{1\leq i<j\leq n}|A_i\cap A_j\cap A_{n+1}|
  +\cdots\\
&\quad+(-1)^{n-1}|A_1\cap A_2\cap\cdots\cap A_n\cap A_{n+1}|.
\end{aligned}
\]
\[
\begin{aligned}
&=\sum_{i=1}^{n+1}|A_i|
  -\sum_{1\leq i<j\leq n+1}|A_i\cap A_j|\\
&\quad+\sum_{1\leq i<j<k\leq n+1}|A_i\cap A_j\cap A_k|
  -\cdots\\
&\quad+(-1)^n|A_1\cap A_2\cap\cdots\cap A_n\cap A_{n+1}|.
\end{aligned}
\]

\bigskip
\noindent\fbox{11}\quad $|M|=m$, $|N|=n$.

\noindent Definim
\[
 s(m,n)=\bigl|\{f:M\longrightarrow N\mid f\text{ surjectivă}\}\bigr|.
\]
\noindent Calculaţi $s(m,n)$. \hfill $s(2,1)=1$.

\noindent\underline{Sol.}
\[
 s(m,n)=0\qquad\text{când }n>m.
\]
\noindent Dacă $n=m$, atunci
\[
 s(m,n)=n!.
\]
\noindent Dacă $n<m$, atunci
\[
 |\operatorname{Hom}(M,N)|=n^m.
\]
\noindent Fie $P=\{f:M\longrightarrow N\mid f\text{ nu este surjectivă}\}$.
Atunci
\[
 s(m,n)=n^m-|P|.
\]
\noindent Scriem
\[
 M=\{a_1,a_2,\ldots,a_m\},\qquad
 N=\{b_1,b_2,\ldots,b_n\}.
\]
\noindent Pentru $i=1,\ldots,n$ notăm
\[
 A_i'=\{f:M\longrightarrow N\mid b_i\notin\operatorname{Im}f\}.
\]
\noindent Rezultă
\[
 |P|=\bigl|A_1'\cup A_2'\cup\cdots\cup A_n'\bigr|,
\]
\[
 |A_i'|=(n-1)^m\quad(\forall i),\qquad
 |A_i'\cap A_j'|=(n-2)^m.
\]
\noindent Aplicând principiul includerii şi excluderii,
\[
\begin{aligned}
 |P|={}&\binom{n}{1}(n-1)^m
       -\binom{n}{2}(n-2)^m\\
      &+\binom{n}{3}(n-3)^m-\cdots
       +(-1)^{n-1}0.
\end{aligned}
\]
.
% Fara preambul si fara intrerupere fortata de pagina.
% NOTE EDITORIALE:
% - Primele formule continua inductia pentru principiul includerii si
%   excluderii din pagina 7. Semnele + din expresia intermediara cu
%   (A_i \cap A_{n+1}) sunt redate ca in manuscris, desi semnul necesar
%   pentru formula finala ar fi minus; nu s-a facut corectarea tacita.
% - In problema 11, manuscrisul foloseste A_i' la definitia multimii
%   de functii care omit b_i, dar omite uneori prima in formulele
%   pentru reuniune si intersectie. In LaTeX se pastreaza A_i' in
%   formulele de numarare pentru a indica aceeasi familie de multimi.
% - La sfarsitul ultimei sume apare un termen taiat / rescris; se reda
%   ultima forma lizibila +(-1)^{n-1} 0, fara a reconstitui termenul taiat.
% - Schemele marginale reprezinta doar exemple de corespondente intre
%   multimi finite; afirmatia explicita s(2,1)=1 este pastrata.

\subsection*{Seminar 2 - Pagina 8}

\noindent\textit{(Continuarea demonstraţiei principiului includerii şi excluderii.)}
\[
\begin{aligned}
&\cdots+(-1)^{n-1}\bigl|A_1\cap\cdots\cap A_n\bigr|\\
&\quad+\bigl|(A_1\cap A_{n+1})\cup(A_2\cap A_{n+1})
    \cup\cdots\cup(A_n\cap A_{n+1})\bigr|\\
&=\sum_{i=1}^{n+1}|A_i|
  -\sum_{1\leq i<j\leq n}|A_i\cap A_j|
  +\cdots+(-1)^{n-1}|A_1\cap\cdots\cap A_n|\\
&\quad+\sum_{i=1}^{n}|A_i\cap A_{n+1}|
  -\sum_{1\leq i<j\leq n}|A_i\cap A_j\cap A_{n+1}|
  +\cdots\\
&\quad+(-1)^{n-1}|A_1\cap A_2\cap\cdots\cap A_n\cap A_{n+1}|.
\end{aligned}
\]
\[
\begin{aligned}
&=\sum_{i=1}^{n+1}|A_i|
  -\sum_{1\leq i<j\leq n+1}|A_i\cap A_j|\\
&\quad+\sum_{1\leq i<j<k\leq n+1}|A_i\cap A_j\cap A_k|
  -\cdots\\
&\quad+(-1)^n|A_1\cap A_2\cap\cdots\cap A_n\cap A_{n+1}|.
\end{aligned}
\]

\bigskip
\noindent\fbox{11}\quad $|M|=m$, $|N|=n$.

\noindent Definim
\[
 s(m,n)=\bigl|\{f:M\longrightarrow N\mid f\text{ surjectivă}\}\bigr|.
\]
\noindent Calculaţi $s(m,n)$. \hfill $s(2,1)=1$.

\noindent\underline{Sol.}
\[
 s(m,n)=0\qquad\text{când }n>m.
\]
\noindent Dacă $n=m$, atunci
\[
 s(m,n)=n!.
\]
\noindent Dacă $n<m$, atunci
\[
 |\operatorname{Hom}(M,N)|=n^m.
\]
\noindent Fie $P=\{f:M\longrightarrow N\mid f\text{ nu este surjectivă}\}$.
Atunci
\[
 s(m,n)=n^m-|P|.
\]
\noindent Scriem
\[
 M=\{a_1,a_2,\ldots,a_m\},\qquad
 N=\{b_1,b_2,\ldots,b_n\}.
\]
\noindent Pentru $i=1,\ldots,n$ notăm
\[
 A_i'=\{f:M\longrightarrow N\mid b_i\notin\operatorname{Im}f\}.
\]
\noindent Rezultă
\[
 |P|=\bigl|A_1'\cup A_2'\cup\cdots\cup A_n'\bigr|,
\]
\[
 |A_i'|=(n-1)^m\quad(\forall i),\qquad
 |A_i'\cap A_j'|=(n-2)^m.
\]
\noindent Aplicând principiul includerii şi excluderii,
\[
\begin{aligned}
 |P|={}&\binom{n}{1}(n-1)^m
       -\binom{n}{2}(n-2)^m\\
      &+\binom{n}{3}(n-3)^m-\cdots
       +(-1)^{n-1}0.
\end{aligned}
\]
.
% Fara preambul si fara intrerupere fortata de pagina.
% NOTE EDITORIALE:
% - Primele formule continua inductia pentru principiul includerii si
%   excluderii din pagina 7. Semnele + din expresia intermediara cu
%   (A_i \cap A_{n+1}) sunt redate ca in manuscris, desi semnul necesar
%   pentru formula finala ar fi minus; nu s-a facut corectarea tacita.
% - In problema 11, manuscrisul foloseste A_i' la definitia multimii
%   de functii care omit b_i, dar omite uneori prima in formulele
%   pentru reuniune si intersectie. In LaTeX se pastreaza A_i' in
%   formulele de numarare pentru a indica aceeasi familie de multimi.
% - La sfarsitul ultimei sume apare un termen taiat / rescris; se reda
%   ultima forma lizibila +(-1)^{n-1} 0, fara a reconstitui termenul taiat.
% - Schemele marginale reprezinta doar exemple de corespondente intre
%   multimi finite; afirmatia explicita s(2,1)=1 este pastrata.

\subsection*{Seminar 2 - Pagina 8}

\noindent\textit{(Continuarea demonstraţiei principiului includerii şi excluderii.)}
\[
\begin{aligned}
&\cdots+(-1)^{n-1}\bigl|A_1\cap\cdots\cap A_n\bigr|\\
&\quad+\bigl|(A_1\cap A_{n+1})\cup(A_2\cap A_{n+1})
    \cup\cdots\cup(A_n\cap A_{n+1})\bigr|\\
&=\sum_{i=1}^{n+1}|A_i|
  -\sum_{1\leq i<j\leq n}|A_i\cap A_j|
  +\cdots+(-1)^{n-1}|A_1\cap\cdots\cap A_n|\\
&\quad+\sum_{i=1}^{n}|A_i\cap A_{n+1}|
  -\sum_{1\leq i<j\leq n}|A_i\cap A_j\cap A_{n+1}|
  +\cdots\\
&\quad+(-1)^{n-1}|A_1\cap A_2\cap\cdots\cap A_n\cap A_{n+1}|.
\end{aligned}
\]
\[
\begin{aligned}
&=\sum_{i=1}^{n+1}|A_i|
  -\sum_{1\leq i<j\leq n+1}|A_i\cap A_j|\\
&\quad+\sum_{1\leq i<j<k\leq n+1}|A_i\cap A_j\cap A_k|
  -\cdots\\
&\quad+(-1)^n|A_1\cap A_2\cap\cdots\cap A_n\cap A_{n+1}|.
\end{aligned}
\]

\bigskip
\noindent\fbox{11}\quad $|M|=m$, $|N|=n$.

\noindent Definim
\[
 s(m,n)=\bigl|\{f:M\longrightarrow N\mid f\text{ surjectivă}\}\bigr|.
\]
\noindent Calculaţi $s(m,n)$. \hfill $s(2,1)=1$.

\noindent\underline{Sol.}
\[
 s(m,n)=0\qquad\text{când }n>m.
\]
\noindent Dacă $n=m$, atunci
\[
 s(m,n)=n!.
\]
\noindent Dacă $n<m$, atunci
\[
 |\operatorname{Hom}(M,N)|=n^m.
\]
\noindent Fie $P=\{f:M\longrightarrow N\mid f\text{ nu este surjectivă}\}$.
Atunci
\[
 s(m,n)=n^m-|P|.
\]
\noindent Scriem
\[
 M=\{a_1,a_2,\ldots,a_m\},\qquad
 N=\{b_1,b_2,\ldots,b_n\}.
\]
\noindent Pentru $i=1,\ldots,n$ notăm
\[
 A_i'=\{f:M\longrightarrow N\mid b_i\notin\operatorname{Im}f\}.
\]
\noindent Rezultă
\[
 |P|=\bigl|A_1'\cup A_2'\cup\cdots\cup A_n'\bigr|,
\]
\[
 |A_i'|=(n-1)^m\quad(\forall i),\qquad
 |A_i'\cap A_j'|=(n-2)^m.
\]
\noindent Aplicând principiul includerii şi excluderii,
\[
\begin{aligned}
 |P|={}&\binom{n}{1}(n-1)^m
       -\binom{n}{2}(n-2)^m\\
      &+\binom{n}{3}(n-3)^m-\cdots
       +(-1)^{n-1}0.
\end{aligned}
\]
